\documentclass[12pt,a4paper]{article}

\usepackage[margin=2.5cm]{geometry}
\usepackage{graphicx}
\usepackage{booktabs}
\usepackage{hyperref}
\usepackage{amsmath}
\usepackage{float}
\usepackage[T1]{fontenc}
\usepackage{lmodern}

\setlength{\parskip}{5pt}
\setlength{\parindent}{0pt}
\linespread{1.2}

\hypersetup{colorlinks=true, linkcolor=blue, urlcolor=blue}

\begin{document}

% ── Title block ──────────────────────────────────────────────────────────────
\begin{titlepage}
    \centering
    \vspace*{2cm}
    {\Large\bfseries Indoor Air Quality in Sport Facilities\par}
    \vspace{0.4cm}
    {\large Interdisciplinary Project Report\par}
    \vspace{2cm}
    \begin{tabular}{ll}
        \textbf{Course:}          & 194.147 -- Interdisciplinary Project in Data Science \\[4pt]
        \textbf{Student:}         & Md.\ Khasrur Rahman (12306556) \\[4pt]
        \textbf{Main Supervisor:} & Ass.-Prof.\ Milica Vujovic, PhD \\[4pt]
        \textbf{Co-Supervisor:}   & Senior Scientist Dr.\ Florina Piroi \\[4pt]
        \textbf{Date:}            & April 2026 \\
    \end{tabular}
    \vfill
\end{titlepage}

% ── 1. Introduction ──────────────────────────────────────────────────────────
\section{Introduction}

Indoor air quality is often overlooked compared to outdoor pollution, but it is particularly important in enclosed sport environments like climbing halls, boulder gyms, and fitness centres. These spaces combine high occupancy, intense physical activity, and sometimes limited ventilation. These conditions can drive up CO$_2$, particulate matter (PM), and temperature to uncomfortable or harmful levels.

The goal of this project is to investigate how key air quality indicators, specifically CO$_2$, particulate matter (PM2.5), temperature, and humidity, vary between sport facility and home environments, and how they relate to occupancy and activity intensity. The dataset comes from a 17-day longitudinal study (June 29 to July 15, 2025) involving 22 participants, each equipped with fixed and portable sensors at both their home and workplace (a sport facility), alongside daily well-being surveys.

The analysis focuses on three questions: (1) How do CO$_2$, PM2.5, temperature and humidity differ between sport facilities and home? (2) Does activity intensity have a measurable effect on CO$_2$? (3) What temporal patterns emerge across the day and the study period?

% ── 2. Research Protocol ─────────────────────────────────────────────────────
\section{Research Protocol}

\subsection{Participants and Study Design}

22 participants (IDs: id088--id108) were monitored at two locations each: their home and their workplace (a sport facility such as a gym or climbing hall). Data collection ran continuously for approximately 17 days, yielding 713 survey observations in total, with 367 recorded at work locations and 260 at home.

\subsection{Sensor Setup, Calibration and Spatial Resolution}

Two sensor types were deployed, providing complementary spatial coverage:

\textbf{AirControl sensors (fixed placement):} One sensor was placed at a fixed indoor position in each participant's home and one at their workplace. These provided continuous fixed-point monitoring of CO$_2$ (ppm), temperature ($^\circ$C), and relative humidity (\%) at high frequency, with timestamps logged in DD/MM/YYYY HH:mm:ss format. Data quality thresholds were applied as a calibration step: CO$_2$ readings outside 300--5000 ppm, temperature outside 10--40$^\circ$C, and humidity outside 10--90\% were treated as implausible and removed.

\textbf{AtmoTube sensors (mobile/portable):} Participants carried these throughout the day. They captured a broader set of indicators: CO$_2$, temperature, humidity, VOC (volatile organic compounds, ppm), air quality score (AQS), and particulate matter PM1, PM2.5, and PM10 ($\mu$g/m$^3$). The portable nature of these sensors provided spatial coverage beyond fixed points, capturing air quality as participants moved through different areas of the facility or their home.

Together, the two sensor types provide both a stable fixed reference (AirControl) and a mobile picture of conditions the participant was actually exposed to (AtmoTube), giving a richer spatial and temporal resolution than either could provide alone.

\subsection{Survey Data Collection}

Participants completed daily surveys on a Likert scale (1--7) across five dimensions: stress, tiredness, productivity, thermal comfort (warmth), and mental demand. These subjective reports were timestamped, allowing them to be matched to the nearest sensor readings.

% ── 3. Methodology ───────────────────────────────────────────────────────────
\section{Methodology}

\subsection{Data Preprocessing}

Raw data was loaded from 321 CSV files across all participants. After applying the sensor thresholds described above, participant IDs were standardised and survey responses outside the valid 1--7 range were removed.

\subsection{Timestamp Alignment}

Sensor readings and survey responses were never recorded at exactly the same moment. To address this, \textbf{linear interpolation} was used to estimate the sensor value at the exact survey timestamp from the two nearest readings on either side:

\[
v_{t} = v_{t_0} + \frac{t - t_0}{t_1 - t_0} \cdot (v_{t_1} - v_{t_0})
\]

This is valid for physical measurements (CO$_2$, temperature, PM2.5) because they change continuously over time. Importantly, this was \textit{not} applied to ordinal survey variables such as stress. A fractional value like ``stress = 3.5'' has no meaningful interpretation since the scale categories are not equally spaced.

\subsection{Correlation Method}

Spearman rank correlation was used throughout rather than Pearson. Survey variables (stress, tiredness, etc.) are ordinal and Pearson assumes equal-interval measurement, which cannot be assumed for Likert-scale data. Spearman converts values to ranks first, making it appropriate for mixed ordinal and continuous variable comparisons.

% ── 4. Results ───────────────────────────────────────────────────────────────
\section{Results}

\subsection{Air Quality Distributions}

Figure~\ref{fig:aq_dist} shows the distributions of the three AirControl sensor variables. CO$_2$ is right-skewed, with the bulk of readings falling below 600 ppm and a long tail reaching above 1400 ppm during peak periods. Temperature clusters around 26--28$^\circ$C and humidity spans a wide range of 15--50\%, reflecting variation across participants' homes and workplaces.

\begin{figure}[H]
    \centering
    \includegraphics[width=0.92\textwidth]{../figures/output2.png}
    \caption{Distribution of AirControl sensor readings across all participants and both locations. CO$_2$ is right-skewed; temperature is approximately normal; humidity shows wide spread.}
    \label{fig:aq_dist}
\end{figure}

\subsection{Work vs.\ Home: Air Quality and Particulate Matter}

Table~\ref{tab:wh} summarises mean values for all key indicators by location type. Figure~\ref{fig:wh_aq} shows the CO$_2$, temperature and humidity distributions visually.

\begin{table}[H]
    \centering
    \caption{Mean air quality indicators by location type.}
    \begin{tabular}{lcc}
        \toprule
        \textbf{Indicator} & \textbf{Work (Sport Facility)} & \textbf{Home} \\
        \midrule
        CO$_2$ (ppm)             & 506.3 & 492.3 \\
        Temperature ($^\circ$C)  & 27.0  & 26.9  \\
        Humidity (\%)            & 30.8  & 31.1  \\
        PM2.5 ($\mu$g/m$^3$)     & 5.13  & 6.16  \\
        \bottomrule
    \end{tabular}
    \label{tab:wh}
\end{table}

\begin{figure}[H]
    \centering
    \includegraphics[width=\textwidth]{../figures/output4.png}
    \caption{CO$_2$, temperature and humidity by location type. CO$_2$ is higher and more variable at work; temperature and humidity are nearly equal across both settings.}
    \label{fig:wh_aq}
\end{figure}

CO$_2$ is consistently higher at work (+14 ppm on average), which is expected given the higher physical activity and occupancy in sport facilities. The sport facility also produces more extreme peaks, with some readings exceeding 1400 ppm, likely during crowded sessions. Temperature and humidity differences between the two settings are negligible.

Particulate matter tells a different story: PM2.5 is slightly \textit{higher at home} (6.16 vs 5.13 $\mu$g/m$^3$). Both values are well below the WHO guideline of 15 $\mu$g/m$^3$, but the higher home value is noteworthy. It may reflect cooking, outdoor air infiltration, or other household sources not present in the sport facility. This suggests that while sport facilities drive CO$_2$ up through occupancy and exertion, home environments contribute more to particulate exposure.

\subsection{Survey Responses by Location}

Figure~\ref{fig:wh_survey} shows well-being scores at work vs.\ home. Stress is slightly higher at work (mean 3.48) compared to home (mean 3.20), but distributions overlap considerably. The other dimensions, including tiredness, productivity, warmth, and mental demand, are similar across both locations.

\begin{figure}[H]
    \centering
    \includegraphics[width=\textwidth]{../figures/output5.png}
    \caption{Survey response distributions by location. Stress is marginally higher at work; other dimensions are largely similar between settings.}
    \label{fig:wh_survey}
\end{figure}

\subsection{Correlations: Air Quality and Well-being}

Figure~\ref{fig:corr} shows the Spearman correlation matrix. Correlations between air quality variables and survey responses are mostly weak (below $|\rho| = 0.15$). The strongest finding within the survey variables is a moderate negative correlation between tiredness and productivity ($\rho \approx -0.41$), which is intuitive. Temperature shows a small positive correlation with tiredness ($\rho \approx 0.31$), suggesting warmer conditions may contribute to fatigue. CO$_2$ alone does not strongly correlate with any single survey dimension.

\begin{figure}[H]
    \centering
    \includegraphics[width=0.72\textwidth]{../figures/output3.png}
    \caption{Spearman correlation matrix between AirControl air quality variables and survey responses. Most air quality and well-being correlations are weak; the strongest pattern is the tiredness--productivity relationship.}
    \label{fig:corr}
\end{figure}

\subsection{Temporal Patterns}

Figure~\ref{fig:temporal} shows daily patterns across the study period. CO$_2$ remains stable with no gradual buildup, suggesting consistent ventilation and stable conditions. Survey responses for stress and tiredness show more day-to-day variation, driven by individual factors rather than environmental drift.

\begin{figure}[H]
    \centering
    \includegraphics[width=\textwidth]{../figures/output6.png}
    \caption{Daily trends across the 17-day study period. CO$_2$ is stable over time; survey responses show more variation driven by individual factors.}
    \label{fig:temporal}
\end{figure}

\subsection{CO$_2$ Thresholds and Well-being}

Grouping observations by CO$_2$ category (Low $<$800, Medium 800--1000, High $>$1000 ppm) reveals a hint that mental demand is slightly higher at elevated CO$_2$ levels (Figure~\ref{fig:co2cat}). However, the Medium and High groups have few observations, so this should be interpreted cautiously.

\begin{figure}[H]
    \centering
    \includegraphics[width=\textwidth]{../figures/output8.png}
    \caption{Survey responses by CO$_2$ category. Most observations fall in the Low category. Mental demand appears slightly elevated in the High CO$_2$ group, but sample sizes limit conclusions.}
    \label{fig:co2cat}
\end{figure}

% ── 5. Discussion and Conclusion ─────────────────────────────────────────────
\section{Discussion and Conclusion}

The results confirm the core expectation: CO$_2$ is higher in sport facilities than at home, driven by occupancy and physical activity. However, the average difference is modest (+14 ppm), and both environments stay mostly within acceptable ranges. The more notable finding is that peak CO$_2$ events (above 800 ppm) occur more frequently in sport facilities, which is a practical concern for facility managers during high-attendance sessions.

Particulate matter shows the opposite pattern, with PM2.5 slightly higher at home. This shows the two environments present different types of air quality risk. Sport facilities raise CO$_2$ through exertion and occupancy, while homes may contribute more to fine particle exposure through domestic sources.

The weak correlations between air quality and subjective well-being were somewhat unexpected but plausible. Most CO$_2$ readings were in the mild range where effects on perception are subtle. The temperature and tiredness correlation ($\rho \approx 0.31$) is the most actionable finding, as warmer sport facility environments may increase perceived fatigue, which has implications for thermal management.

Methodologically, the two key decisions, using Spearman correlation for ordinal survey data and using linear interpolation rather than window averaging for timestamp alignment, were important for producing statistically valid results. Future work could improve further by collecting occupancy counts directly, increasing sensor sampling frequency, and extending the study period to capture seasonal variation.

In summary, this project provided a data-driven picture of indoor air quality in sport facilities. CO$_2$ and PM2.5 vary in opposite directions between work and home, temporal patterns track occupancy hours, and subjective well-being shows only weak ties to short-term air quality fluctuations at these moderate exposure levels.

\end{document}
